% ============================================================
%  Computational Hardware Effects on Autonomous Vehicle Control:
%  A SHARC–CARLA Co-simulation Framework
%  IEEE CDC submission — double-column, 10pt
%  Authors: [PLACEHOLDER — fill in before submission]
% ============================================================
\documentclass[conference]{IEEEtran}

% ---- Packages -----------------------------------------------
\usepackage{amsmath,amssymb,amsthm}
\usepackage{bm}          % bold math
\usepackage{graphicx}
\usepackage{booktabs}
\usepackage{multirow}
\usepackage{xcolor}
\usepackage{hyperref}
\usepackage{cite}
\usepackage{algorithm}
\usepackage{algpseudocode}
\usepackage[T1]{fontenc}

% ---- Macros --------------------------------------------------
\newcommand{\R}{\mathbb{R}}
\newcommand{\N}{\mathbb{N}}
\newcommand{\ZOH}{\mathrm{ZOH}}
\newcommand{\RMSE}{\mathrm{RMSE}}
\newcommand{\diag}{\mathrm{diag}}
\newcommand{\calX}{\mathcal{X}}
\newcommand{\calU}{\mathcal{U}}
\newcommand{\TODO}[1]{\textcolor{red}{\textbf{[TODO: #1]}}}

% ---- Theorem environments -----------------------------------
\newtheorem{definition}{Definition}
\newtheorem{remark}{Remark}

% =============================================================
\begin{document}

\title{Computational Hardware Effects on Autonomous Vehicle
  Control: A SHARC--CARLA Co-simulation Framework}

\author{\IEEEauthorblockN{Author~1, Author~2, Author~3}
  \IEEEauthorblockA{Affiliation Placeholder\\
    \texttt{email@placeholder.edu}}
}

\maketitle

% =============================================================
\begin{abstract}
Autonomous vehicles (AVs) execute safety-critical control algorithms—PID
regulators, linear MPC, nonlinear MPC—on a wide spectrum of embedded
processors, from high-performance automotive SoCs to constrained
microcontrollers.  The computation time $\tau_k$ incurred by these
algorithms is not fixed; it depends intimately on the underlying
microarchitecture (cache hierarchy, clock frequency, branch predictor).
A delayed control input is effectively a stale input: the plant evolves
for an entire sample period on the wrong command.  In safety-critical
scenarios—emergency braking, intersection crossing, high-speed path
following—even a single missed computational deadline can be catastrophic.
This paper presents a unified co-simulation framework that couples
\emph{SHARC} (Simulator for Hardware Architecture and Real-time
Control)~\cite{wintz2025sharc} with the \emph{CARLA} high-fidelity
urban driving simulator.  We model CARLA as a concrete \emph{dynamics
class} inside SHARC, replacing its ODE integrator with cycle-accurate
physics ticks.  We then deploy three controllers of increasing
computational complexity (PID, LMPC, NLMPC) across multiple hardware
configurations and AV-relevant driving scenarios, measuring trajectory
tracking error, deadline miss rate, and safety event frequency.  Our
results reveal a consistent \emph{hardware--performance cliff}: small
reductions in processor clock speed or cache size push the computational
delay beyond the sample period, triggering cascading control degradation.
This framework enables quantitative hardware--software co-design for AVs
before silicon is committed.
\end{abstract}

\begin{IEEEkeywords}
Autonomous driving, hardware--software co-design, model predictive
control, computational delay, CARLA, microarchitecture simulation.
\end{IEEEkeywords}

% =============================================================
\section{Introduction}\label{sec:intro}

Modern autonomous vehicles are cyber-physical systems (CPSs) in the
strictest sense: their safety properties emerge from the \emph{joint}
behavior of physics, control software, and computational hardware.  Yet
the dominant design practice treats these three as separable concerns.
Control engineers design algorithms assuming zero-delay actuation;
hardware engineers provision processors for average-case workloads;
systems integrators discover timing violations late—often on physical
prototypes.

The root cause of this disconnect is the lack of a simulation tool that
simultaneously models (a) realistic vehicle dynamics, (b) the
computation time of the control algorithm on a specific processor, and
(c) the causal effect of that delay on the closed-loop trajectory.
Existing alternatives fall short.  Hardware-in-the-loop (HIL) systems
use real hardware, making them inflexible and expensive for design-space
exploration~\cite{mihalic2022hil}.  Standard co-simulators lack
cycle-accurate timing~\cite{torngren2006cosim}.  CARLA alone models
physics beautifully but treats control computation as instantaneous.

We close this gap by integrating two state-of-the-art simulators:
\textbf{SHARC}~\cite{wintz2025sharc}, which uses the Scarab
microarchitecture simulator~\cite{scarab} to measure controller
execution time at the instruction level, and \textbf{CARLA}~\cite{dosovitskiy2017carla},
which provides a photorealistic, physics-faithful urban driving
environment.  The key contribution is a \texttt{CarlaDynamics}
interface class that exposes CARLA's synchronous stepping API as a
SHARC \texttt{Dynamics} object, so that every SHARC simulation loop
tick advances CARLA by exactly one sample period, while Scarab
concurrently measures how long the controller binary took to run on
the target chip.

\smallskip\noindent\textbf{Contributions.}
\begin{enumerate}
  \item \textbf{SHARC--CARLA integration}: a reusable \texttt{CarlaDynamics}
    class that bridges cycle-accurate hardware simulation with
    high-fidelity autonomous-driving physics.
  \item \textbf{Controller library}: PID, LMPC, and NLMPC controllers
    formulated uniformly within the SHARC pipe-based communication
    protocol, enabling apples-to-apples hardware comparison.
  \item \textbf{Hardware--software co-design analysis}: a systematic
    study of how compute budget (clock speed, cache) interacts with
    controller type and scenario complexity, quantified by tracking
    error, deadline miss rate, and safety event frequency.
\end{enumerate}

% =============================================================
\section{Background: The SHARC Framework}\label{sec:sharc}

We summarize the formal model of SHARC~\cite{wintz2025sharc}, which
this work extends to autonomous driving.

\subsection{Plant Model}

A physical plant is modeled in continuous time as
\begin{equation}\label{eq:ct}
  \dot{x} = \tilde{f}(t, x, u, w), \qquad y = h(x, u, w),
\end{equation}
where $x \in \R^{n_x}$ is the state, $u \in \R^{n_u}$ the control
input, $w \in \R^{n_w}$ an exogenous (disturbance/reference) input, and
$y \in \R^{n_y}$ the output.  SHARC discretizes~\eqref{eq:ct} on a
uniform grid $t_k := kT$ with sample period $T > 0$:
\begin{equation}\label{eq:disc}
  x_{k+1} = f(t_k, x_k, u_k, w(t_k)),
\end{equation}
where $f$ is obtained by numerical integration of $\tilde{f}$ over
$[t_k, t_{k+1}]$ under zero-order hold (ZOH) on $u$.  Users may
alternatively supply $f$ directly (hybrid systems, stochastic models).

\subsection{Computation-Delay Model}\label{sec:delay}

In an ideal controller, $u_k = g(t_k, y_k)$ is available
instantaneously.  In reality, computing $g(t_k, y_k)$ requires
$\tau_k > 0$ seconds.  Let $\tilde{t}_{k+1} := t_k + \tau_k$ be the
time at which the new control is ready.  SHARC models the discrete
closed-loop as:
\begin{equation}\label{eq:zoh_delay}
  u(t) = \begin{cases}
    u_k & t \in [t_k,\; t_{k+1}) \text{ if } \tau_k \geq T, \\
    u_k & t \in [t_k,\; \tilde{t}_{k+1}) \text{ if } \tau_k < T, \\
    \tilde{u}_{k+1} & t \in [\tilde{t}_{k+1},\; t_{k+1}) \text{ if }
    \tau_k < T,
  \end{cases}
\end{equation}
where $\tilde{u}_{k+1} := g(t_k, y_k)$ is the freshly computed control.
When $\tau_k \geq T$ (a \emph{deadline miss}), the stale input $u_k$
is held for the entire next period.  Cascading misses cause the system
to accumulate error, often destabilizing it.

\begin{definition}[Deadline Miss]\label{def:miss}
  At time step $k$, a \emph{deadline miss} occurs if $\tau_k \geq T$.
  The \emph{miss rate} over a horizon of $K$ steps is
  $\rho := \frac{1}{K}\sum_{k=0}^{K-1}\mathbf{1}[\tau_k \geq T]$.
\end{definition}

\subsection{Computation Time via Scarab}

The delay $\tau_k$ is not assumed; it is \emph{measured} by running the
controller binary inside the Scarab microarchitecture
simulator~\cite{scarab}.  Scarab is cycle-accurate: it models branch
predictors, cache hierarchies, and instruction pipelines, translating
them into a simulated execution time in femtoseconds.  Two modes exist:
\emph{execution-driven} (Scarab instruments the running binary via
processor-trace markers) and \emph{DynamoRIO-based} (memory traces are
captured offline and replayed in parallel batches).  Hardware is
specified via \texttt{PARAMS} files: clock frequency, L1/L2 cache
sizes, and memory latency.  This makes it straightforward to evaluate a
controller on, say, an ARM Cortex-A76 versus a Cortex-M55 without
owning either chip.

% =============================================================
\section{CARLA as a SHARC Dynamics Class}\label{sec:carla_dyn}

\subsection{Bridging CARLA and SHARC}

CARLA is a synchronous open-source vehicle simulator built on Unreal
Engine 4.  In synchronous mode, time advances only when the Python
client calls \texttt{world.tick()}, making it amenable to tight
integration with an external timing model.

We implement a \texttt{CarlaDynamics} class that inherits from
SHARC's abstract \texttt{Dynamics} base class.  The interface contract
requires two methods:
\begin{itemize}
  \item \texttt{setup\_system()}: called once before the simulation loop;
    connects to the CARLA server, sets $T_s =$ \texttt{fixed\_delta\_seconds},
    spawns the ego vehicle, and initializes sensors.
  \item \texttt{evolve\_state}$(t_0, x_0, u, w, t_f)$: applies control
    $u$ to the ego vehicle, calls \texttt{world.tick()} for
    $\lfloor(t_f - t_0)/T_s\rfloor$ steps, and reads back
    the new state $x_{k+1}$ from the CARLA actor API.
\end{itemize}

Critically, CARLA is put into \emph{synchronous mode with fixed
  delta}: the internal physics engine advances by exactly $T$ per tick,
ensuring that SHARC's discrete time index $k$ and CARLA's physics time
remain aligned throughout the simulation.

\subsection{State, Control, and Exogenous Input}

For the ego vehicle, we define:
\begin{align}
  x_k &:= \bigl[p_x,\; p_y,\; \psi,\; v,\; w_x^{\mathrm{nxt}},\;
    w_y^{\mathrm{nxt}}\bigr]^\top \in \R^6, \label{eq:state}\\
  u_k &:= \bigl[\alpha,\; \delta,\; \beta\bigr]^\top \in
    [0, 1] \times [-1, 1] \times [0, 1], \label{eq:control}\\
  w_k &:= \bigl[w_x^{\mathrm{nxt}},\; w_y^{\mathrm{nxt}}\bigr]^\top
    \in \R^2, \label{eq:exo}
\end{align}
where $(p_x, p_y)$ is the ego position, $\psi$ the yaw angle,
$v$ the longitudinal speed in km/h, $(w_x^{\mathrm{nxt}},
w_y^{\mathrm{nxt}})$ the coordinates of the next road waypoint
obtained from the CARLA HD map (2\,m ahead), and $(\alpha, \delta,
\beta)$ are throttle, steering, and brake commands.  The output is
$y_k := [p_x, p_y]^\top \in \R^2$.

\begin{remark}
  Embedding the next waypoint in both state and exogenous input is deliberate:
  the C++ controller receives $w_k$ via the SHARC pipe interface without
  a separate communication channel, and the Python plant runner can
  independently query it from the CARLA map for state bookkeeping.
\end{remark}

The transition $x_{k+1} = F_{\mathrm{CARLA}}(x_k, u_k, w_k)$ replaces
the ODE integrator in~\eqref{eq:disc}: CARLA's UE4 physics engine acts
as an implicit, high-fidelity integrator that internally models tire
forces, aerodynamic drag, engine torque curves, and road geometry.

% =============================================================
\section{Controller Formulations}\label{sec:controllers}

All three controllers are compiled as standalone C++ executables and
communicate with the SHARC plant runner via named POSIX pipes (FIFOs).
At each step $k$, the Python plant runner writes $(k, t_k, x_k, w_k)$
to input pipes; the C++ binary reads them, computes $u_k$, and writes
it to the output pipe.  Scarab instruments this computation window to
obtain $\tau_k$.

\subsection{PID Controller}

We decompose vehicle control into two independent PID loops.

\noindent\textbf{Longitudinal (speed) control.}
\begin{equation}\label{eq:pid_lon}
  e_v^k := v^* - v_k, \qquad
  \sigma_k = K_P^{\ell} e_v^k + K_D^{\ell}\,\dot{e}_v^k +
  K_I^{\ell}\int_0^{t_k} e_v \,d\tau.
\end{equation}
If $\sigma_k \geq 0$, set $\alpha_k = \min(\sigma_k, \alpha_{\max})$,
$\beta_k = 0$; otherwise $\alpha_k = 0$, $\beta_k = \min(-\sigma_k,
\beta_{\max})$.

\noindent\textbf{Lateral (steering) control.}  The heading error to the
next waypoint is
\begin{equation}\label{eq:heading_err}
  e_\psi^k := \mathrm{atan2}(w_y^k - p_y^k,\;w_x^k - p_x^k) - \psi_k,
\end{equation}
and the steering command is
\begin{equation}\label{eq:pid_lat}
  \delta_k = K_P^{r} e_\psi^k + K_D^{r}\,\dot{e}_\psi^k +
  K_I^{r}\int_0^{t_k} e_\psi \,d\tau,
\end{equation}
clipped to $[-\delta_{\max}, \delta_{\max}]$.  Derivative and integral
terms are computed online using a sliding buffer of the last
$N_{\mathrm{buf}} = 10$ samples.

\subsection{Linear MPC (LMPC)}

We linearize the CARLA vehicle dynamics around a reference trajectory
$(\bar{x}_k, \bar{u}_k)$ to obtain the incremental model
\begin{equation}\label{eq:lmpc_model}
  \Delta x_{k+1} = A_k \Delta x_k + B_k \Delta u_k,
\end{equation}
where $\Delta x_k := x_k - \bar{x}_k$, $\Delta u_k := u_k -
\bar{u}_k$, and $(A_k, B_k)$ are the Jacobians of
$F_{\mathrm{CARLA}}$ evaluated at the reference point.

At each sample time $t_{k_0}$, LMPC solves the convex QP:
\begin{align}
  \min_{\{u_k\}_{k=k_0}^{k_0+N_c-1}} \quad&
    \sum_{k=k_0}^{k_0+N_p-1}\!\!\|\Delta x_{k|k_0}\|_Q^2
    + \|\Delta u_{k|k_0}\|_R^2 \label{eq:lmpc_cost}\\
  \mathrm{s.t.}\quad & \Delta x_{k+1|k_0} = A_k\Delta x_{k|k_0} +
    B_k \Delta u_{k|k_0}, \nonumber\\
  & \Delta x_{k_0|k_0} = x_{k_0} - \bar{x}_{k_0},\nonumber\\
  & u_k \in \mathcal{U},\quad x_k \in \mathcal{X}, \nonumber
\end{align}
where $N_p$ is the prediction horizon, $N_c \leq N_p$ the control
horizon, $Q \succeq 0$ and $R \succ 0$ are state and input cost
matrices, $\mathcal{U}$ encodes actuator limits, and $\mathcal{X}$
encodes road/safety constraints (lane boundaries, speed limits).  Only
$u_{k_0|k_0}$ is applied (receding horizon).

\subsection{Nonlinear MPC (NLMPC)}

NLMPC retains the full nonlinear transition:
\begin{align}
  \min_{\{u_k\}_{k=k_0}^{k_0+N_c-1}} \quad&
    \sum_{k=k_0}^{k_0+N_p-1} C(x_{k|k_0}, u_{k|k_0})
    \label{eq:nlmpc_cost}\\
  \mathrm{s.t.}\quad & x_{k+1|k_0} =
    F_{\mathrm{model}}(x_{k|k_0}, u_{k|k_0}), \nonumber\\
  & x_{k_0|k_0} = \hat{x}_{k_0},\nonumber\\
  & u_k \in \mathcal{U},\quad x_k \in \mathcal{X}. \nonumber
\end{align}
Here $F_{\mathrm{model}}$ is a \emph{kinematic bicycle model} used
as the internal prediction model (not CARLA itself), providing a
computationally tractable nonlinear approximation.  The cost
$C(x_k, u_k) = \|x_k - x_k^{\mathrm{ref}}\|_Q^2 +
\|u_k\|_R^2$ is a reference-tracking quadratic.  The optimization
is solved by Sequential Least Squares QP (SLSQP)~\cite{kraft1988slsqp}
via the \texttt{libmpc++} library~\cite{piccinelli_libmpc}.

\begin{remark}[Computational complexity ordering]
  PID has $O(1)$ arithmetic operations per step; LMPC solves a convex QP
  of size $O(N_p n_x + N_c n_u)$; NLMPC iteratively solves a nonconvex
  NLP.  This ordering ($\tau_k^{\mathrm{PID}} \ll
  \tau_k^{\mathrm{LMPC}} \ll \tau_k^{\mathrm{NLMPC}}$) is central to
  our hardware sensitivity analysis.
\end{remark}

% =============================================================
\section{Experimental Setup}\label{sec:experiments}

\subsection{Driving Scenarios}

Table~\ref{tab:scenarios} summarizes the six scenarios drawn from
CARLA's built-in maps and town configurations.  Scenarios are ordered
by \emph{safety criticality}: the higher the row, the less margin is
available for delayed actuation.

\begin{table}[h]
  \centering
  \caption{Driving scenarios used for evaluation.}
  \label{tab:scenarios}
  \renewcommand{\arraystretch}{1.2}
  \footnotesize
  \begin{tabular}{llp{2.6cm}c}
    \toprule
    \textbf{Scenario} & \textbf{Map} & \textbf{Key challenge}
      & \textbf{$T$ (s)}\\
    \midrule
    Straight-road cruising     & Town04 & Speed tracking        & 0.25 \\
    Highway lane following     & Town04 & Lateral precision     & 0.25 \\
    Adaptive cruise control    & Town04 & Lead-car gap          & 0.20 \\
    Intersection crossing      & Town03 & Timing, collision av. & 0.15 \\
    Emergency braking          & Town01 & Min. stopping dist.   & 0.10 \\
    High-speed race lap        & Town04 & Stability at speed    & 0.10 \\
    \bottomrule
  \end{tabular}
\end{table}

For each scenario a reference trajectory $\{(\bar{x}_k,
\bar{u}_k)\}_{k=0}^K$ is pre-computed using CARLA's built-in planner
and stored offline.

\subsection{Hardware Configurations}

Each controller binary is simulated under three hardware configurations
expressed as Scarab \texttt{PARAMS} files (Table~\ref{tab:hw}).
We additionally sweep the clock frequency from a \emph{baseline}
(sufficient for all controllers) down to a \emph{constrained} regime
to identify each controller's \emph{minimum viable clock speed}.

\begin{table}[h]
  \centering
  \caption{Hardware configurations used in Scarab.}
  \label{tab:hw}
  \renewcommand{\arraystretch}{1.2}
  \footnotesize
  \resizebox{\columnwidth}{!}{%
  \begin{tabular}{llll}
    \toprule
    \textbf{Config} & \textbf{Clock} & \textbf{L1-D cache}
      & \textbf{Archetype}\\
    \midrule
    HW-High  & \TODO{X} GHz & \TODO{X} kB & High-perf.\ automotive SoC \\
    HW-Mid   & \TODO{X} GHz & \TODO{X} kB & Mid-range ARM Cortex-A  \\
    HW-Low   & \TODO{X} MHz & \TODO{X} kB & Microcontroller-class    \\
    \bottomrule
  \end{tabular}}%
\end{table}

\subsection{Performance Metrics}

Let $\{x_k\}_{k=0}^K$ be the simulated trajectory and
$\{\bar{x}_k\}_{k=0}^K$ the reference.

\noindent\textbf{Position RMSE:}
\begin{equation}\label{eq:rmse}
  \RMSE_{\mathrm{pos}} :=
  \sqrt{\frac{1}{K}\sum_{k=0}^{K-1}
    \bigl[(p_x^k - \bar{p}_x^k)^2 + (p_y^k - \bar{p}_y^k)^2\bigr]}.
\end{equation}

\noindent\textbf{Speed RMSE:}
$\RMSE_v := \sqrt{\frac{1}{K}\sum_{k=0}^{K-1}(v_k - v^*)^2}$.

\noindent\textbf{Deadline miss rate} $\rho$ (Definition~\ref{def:miss}).

\noindent\textbf{Mean computation delay:}
$\bar{\tau} := \frac{1}{K}\sum_{k=0}^{K-1} \tau_k$, with distribution
$p(\tau_k)$ visualized as a box plot.

\noindent\textbf{Safety events:} number of (a) collisions with the
environment, (b) lane-boundary violations ($|e_y| > 1.5$\,m), and
(c) speed-limit exceedances ($v > 1.2\,v^*$) per episode.

\subsection{Implementation Details}

The full framework runs inside a single Docker container
(\texttt{carla-sharc}) built on \texttt{carlasim/carla:0.9.16}
(Ubuntu 20.04, CUDA 11, Python 3.10 via Conda) with Scarab built from
commit \texttt{3b38da0}.  The CARLA server is launched in synchronous,
\texttt{-RenderOffScreen} mode for headless experimentation.  All
experiments use the execution-driven Scarab mode (no DynamoRIO) for
simplicity; the parallel batch mode can be enabled for NLMPC to reduce
wall-clock time as in~\cite{wintz2025sharc}.  All controller source
code, PARAMS files, and docker configuration are available at
\texttt{\TODO{github.com/PLACEHOLDER/sharc-carla}}.

% =============================================================
\section{Results}\label{sec:results}

\subsection{Computation Delay Distributions}

\TODO{Figure: Box plots of $\tau_k$ distributions for PID, LMPC, NLMPC
across HW-High, HW-Mid, HW-Low on the lane-following scenario.  Annotate
with horizontal line at $T_s = 0.25$\,s to show deadline-miss threshold.
Expected finding: PID always below threshold; NLMPC exceeds on HW-Low.}

\begin{table}[h]
  \centering
  \caption{Mean computation delay $\bar{\tau}$ and miss rate $\rho$
    (lane-following, $T=0.25$\,s).}
  \label{tab:delays}
  \footnotesize
  \resizebox{\columnwidth}{!}{%
  \begin{tabular}{lcccccc}
    \toprule
    & \multicolumn{2}{c}{\textbf{PID}}
    & \multicolumn{2}{c}{\textbf{LMPC}}
    & \multicolumn{2}{c}{\textbf{NLMPC}}\\
    \cmidrule(lr){2-3}\cmidrule(lr){4-5}\cmidrule(lr){6-7}
    \textbf{HW} & $\bar\tau$ & $\rho$ & $\bar\tau$ & $\rho$
      & $\bar\tau$ & $\rho$ \\
    \midrule
    High & \TODO{} & \TODO{} & \TODO{} & \TODO{} & \TODO{} & \TODO{} \\
    Mid  & \TODO{} & \TODO{} & \TODO{} & \TODO{} & \TODO{} & \TODO{} \\
    Low  & \TODO{} & \TODO{} & \TODO{} & \TODO{} & \TODO{} & \TODO{} \\
    \bottomrule
  \end{tabular}}% end resizebox
\end{table}

\subsection{Tracking Performance vs.\ Hardware}

\TODO{Figure: RMSE$_\text{pos}$ vs.\ clock frequency sweep for
NLMPC and LMPC on the intersection and emergency-braking scenarios.
Expected finding: step-function degradation at a critical clock—the
``hardware--performance cliff.''  PID flat across all frequencies,
confirming computational simplicity as robustness.}

\begin{table}[h]
  \centering
  \caption{Position RMSE (m) across scenarios and hardware configs.}
  \label{tab:rmse}
  \footnotesize
  \resizebox{\columnwidth}{!}{%
  \begin{tabular}{lcccccc}
    \toprule
    & \multicolumn{3}{c}{\textbf{LMPC}}
    & \multicolumn{3}{c}{\textbf{NLMPC}}\\
    \cmidrule(lr){2-4}\cmidrule(lr){5-7}
    \textbf{Scenario} & High & Mid & Low & High & Mid & Low \\
    \midrule
    Lane following    & \TODO{} & \TODO{} & \TODO{} & \TODO{} & \TODO{} & \TODO{} \\
    ACC               & \TODO{} & \TODO{} & \TODO{} & \TODO{} & \TODO{} & \TODO{} \\
    Intersection      & \TODO{} & \TODO{} & \TODO{} & \TODO{} & \TODO{} & \TODO{} \\
    Emerg.\ braking   & \TODO{} & \TODO{} & \TODO{} & \TODO{} & \TODO{} & \TODO{} \\
    High-speed race   & \TODO{} & \TODO{} & \TODO{} & \TODO{} & \TODO{} & \TODO{} \\
    \bottomrule
  \end{tabular}}% end resizebox
\end{table}

\subsection{Safety Events}

\TODO{Table: Collision count, lane-departure count, speed-limit
exceedance count for each controller$\times$HW pair across all
scenarios.  Expected finding: NLMPC on HW-Low shows highest collision
rate in emergency braking and intersection scenarios; PID shows most
lane departures in high-speed race due to limited optimality.}

\subsection{Hardware--Software Co-design Analysis}

\TODO{Figure: Pareto frontier plots: RMSE$_\text{pos}$ vs.\ clock
speed for LMPC and NLMPC with varying prediction horizon $N_p \in
\{5, 10, 20\}$.  Each $(N_p, f_\text{clk})$ pair is a design point.
Expected finding: smaller $N_p$ reduces computation time, lowers
$\rho$, but increases tracking error; the Pareto frontier identifies
optimal $(N_p, f_\text{clk})$ tradeoffs for each AV scenario.}

This analysis operationalizes the co-design objective: rather than
choosing the largest $N_p$ that a given processor can \emph{nominally}
run, the SHARC--CARLA framework reveals the largest $N_p$ the processor
can run \emph{within the sample period} of a specific driving scenario,
accounting for worst-case instruction-level timing.

% =============================================================
\section{Discussion}\label{sec:discussion}

\noindent\textbf{The hardware--performance cliff.}
A recurring observation in SHARC-based experiments (original paper:
inverted pendulum at 600\,MHz vs.\ 625\,MHz~\cite{wintz2025sharc}) is
that degradation is not graceful.  When $\tau_k$ crosses $T$ once, the
resulting suboptimal control input perturbs the state, which increases
the optimization difficulty at the next step, which in turn increases
$\tau_{k+1}$—a positive feedback loop.  In autonomous driving, this
manifests as spinning controllers in emergency braking: exactly the
scenario where reliability matters most.  PID, by contrast, has a
deadline miss rate near zero across all hardware configurations, making
it a robust fallback despite its suboptimal tracking.

\noindent\textbf{MPC horizon as a co-design variable.}
The classical MPC design question ``how large should $N_p$ be?'' is
answered in closed-loop terms by our Pareto analysis.  On constrained
hardware, halving $N_p$ can eliminate all deadline misses while
incurring only a modest RMSE increase—a trade that the designer cannot
evaluate without SHARC.

\noindent\textbf{Limitations.}
CarlaDynamics uses CARLA's Python API, introducing API-call latency
that is excluded from Scarab's measurement window.  The internal CARLA
physics engine is a UE4 rigid-body model, not a certified vehicle
model; results may not directly transfer to production vehicles.
Future work: (1)~extend \texttt{CarlaDynamics} to multi-agent traffic,
enabling analysis of inter-vehicle communication delays; (2)~introduce
a formal safety guarantee based on control barrier functions compatible
with SHARC's delay model.

% =============================================================
\section{Conclusion}\label{sec:conclusion}

We presented a SHARC--CARLA co-simulation framework for studying the
joint effect of computational hardware and control algorithm design on
autonomous vehicle safety.  By implementing CARLA as a native SHARC
dynamics class and deploying PID, LMPC, and NLMPC controllers under
cycle-accurate microarchitecture simulation, our framework enables
quantitative hardware--software co-design before any physical hardware
is committed.  Results \TODO{(to be completed)} are expected to
demonstrate a sharp hardware--performance cliff for MPC-class
controllers, Pareto-optimal $(N_p, f_\text{clk})$ design surfaces, and
the complementary resilience of PID under extreme hardware constraints.
We believe this methodology generalizes beyond AVs to any CPS where
computation-intensive safety-critical algorithms need to be deployed on
resource-constrained embedded processors.

% =============================================================
\bibliographystyle{IEEEtran}
\bibliography{references}

\end{document}
